\section{The original second layer and its integral ramified model}
\label{sec:second-layer-global}

\status{RSL1--RSL6 identify the exact mass, volume and finite-character
inputs on the original records.  RSL7--RSL10 prove the integral and
isometric transport from the actual level-zero forms to the augmentation
model; this supplies the previously retained transport hypothesis in
LC7.  RSL11--RSL12 give a smaller sufficient condition for actual
\texttt{NuFullF} and \texttt{MuRank}.  RSL13 records the exact remaining
global slope, operator and assembly obligations.  A theorem about a
supplied second layer is not being counted as construction of the common
family.  All references to an unresolved obligation concern this human
reconstruction, under the standing assumption that the pinned formal
proof is correct.}

Throughout, $D$ is the original \texttt{D7} record, $E/F$ its CM
extension, $K=D.\texttt{what}_2.\texttt{LL}$ its coefficient field, and
$s_K:K\hookrightarrow\mathbb C$ its specified embedding.  The
normalizations of mass, Haar measure and squared Hermitian slope are
those already fixed in MF, GV and LC7.  In particular, a new numerical
function called a volume, an abstractly isomorphic form space, or a
finite-index sublattice cannot replace a field of $D$ in the arguments
below.

\subsection{Mass and volume on the very same Hermitian spaces}

% BEGIN RSL1
\paragraph{RSL1: the literal genus-mass input in every rank.}
\label{claim:RSL1}
For every original $D$ and every integer $n\ge0$, the following
statement is established:
\[
 \begin{gathered}
 W:\texttt{HSpace}(E),\quad W.\texttt{SigOK},\quad
 W.\mathrm{sig}(\tau)=(W.\mathrm{rank},0)\quad\hbox{for every }\tau,
 \quad W.\mathrm{rank}=n\quad\Longrightarrow\\
 \operatorname{mass}(W.h_T,W.L_0)=\operatorname{massI}(D,n).
 \end{gathered}                                                   \tag{RSL.1}
\]
This is exactly \texttt{GenusMassAt D n}; in particular it supplies
the universally quantified input \texttt{hG} in \texttt{link\_gl}.

Here is the full transport, including the identities of the lattices.
The hypotheses \texttt{SigOK} and the CM-unramified field of $D$
place the actual pair $(W.h_T,W.L_0)$ in the scope of GV8--GV11.
GV11 identifies the normalized adelic quotient with the genus of this
very lattice: a finite adelic representative $g$ gives
$W\cap g_f\widehat {W.L_0}$, and its rational stabilizer is its full
lattice automorphism group.  Consequently its quotient measure is
the sum of the reciprocal orders of those automorphism groups, which
is precisely the left side of (RSL.1).  It is not the mass of a
replacement lattice.  For $n>0$, GV9--GV10 give that quotient the
value $M_n$ of (GV.21).  Apply the identical calculation to the
standard positive pair $(E^n,I,\mathcal O_E^n)$.  The definition of
\texttt{massI} is its actual genus mass, so it too is $M_n$.
For $n=0$, both spaces and lattices are zero, their automorphism group
is the one-element group, and both masses are one.  This proves
(RSL.1) in every rank.

The retained published input is exactly GV-TAM, with the measure
normalization already checked in GV9--GV10.  The use of that theorem
does not require a global isometry between $W$ and the standard
positive space.  Nor does (RSL.1) assume \texttt{VolTie} or change
$D.\mathrm{vol}$: this is a statement about the actual two genus
masses.
% END RSL1

% BEGIN RSL2
\paragraph{RSL2: the literal reciprocal volume ratio.}
\label{claim:RSL2}
For every original $D$, GV14--GV15 give
\[
 \texttt{VolTie}(D)\ \Longrightarrow\ \texttt{VolumeRatio}_2(D),
                                                                  \tag{RSL.2}
\]
which is exactly \texttt{MassRatioUnderVolume D}, the input
\texttt{hU} of \texttt{link\_gl}.  We check both the witness and the
normalization rather than identifying two statements by their names.

Let $W$ have constant signature $(P,Q)$ and rank $n\ge3$, and let
$\epsilon\in\{1,-1\}$ be a sign represented by that signature.
GV14, with its retained Hasse--Minkowski input MF-HM, constructs a
rational $x\in W$ with $h_T(x,x)=\epsilon$.  It constructs a
self-dual lattice on the actual rational complement $W'=x^\perp$
and an isometric inclusion $\iota:W'\hookrightarrow W$ with range
$\ker h_T(x,-)$.  The signatures of $W'$ are $(P-1,Q)$ or
$(P,Q-1)$ according to the chosen sign.  These are exactly the
vector, complement, inclusion, range and signature witnesses that
occur in \texttt{VolumeRatio}\textsubscript{2}; none is reselected
when we pass to numerical volumes.

Choose normalized Haar measures and fundamental domains for $W$ and
$W'$ using the existence clauses of the original \texttt{VolTie}.
GV15 applies to every such pair of measures and domains.  The other
clauses of \texttt{VolTie} identify their measures with
$\operatorname{ofReal}(D.\mathrm{vol}(W))$ and
$\operatorname{ofReal}(D.\mathrm{vol}(W'))$.  The original
$D.\texttt{vol\_pos}$ fields justify taking real values.  Thus, with
$v=D.\mathrm{host}.v>0$, $f=[F:\mathbb Q]$ and
$L_n=\mathrm{lValEta}(E,n)>0$, GV15 yields
\[
 \frac{D.\mathrm{vol}(W)}{D.\mathrm{vol}(W')}
 =\frac{\big((n-1)!(v/\pi)^n\big)^f}{\sqrt{(2v)^f}}L_n.
                                                                  \tag{RSL.3}
\]
Taking reciprocals, using positive real powers, gives
\[
 \boxed{\quad
 \frac{D.\mathrm{vol}(W')}{D.\mathrm{vol}(W)}
   =\left(\frac{(2\pi)^n}{(n-1)!}\right)^f
        \big((2v)^f\big)^{1/2-n}L_n^{-1}.
 \quad}                                                          \tag{RSL.4}
\]
Indeed the ratio of the right sides of (RSL.4) and the reciprocal
of (RSL.3) is $2^{nf}(2v)^{-nf}v^{nf}=1$.
This is the precise coefficient in the source, including the
orientation of the ratio and its half power.  The proof works for
every $D$ and every admissible $W$, so the universal quantifier in
\texttt{hU} has not been weakened to a chosen family member.
Its retained inputs are GV-TAM and MF-HM, both already explicit in
GV14--GV15.
% END RSL2

% BEGIN RSL3
\paragraph{RSL3: real-place and base-volume transport without retieing.}
\label{claim:RSL3}
For the same $D$,
\[
 \texttt{RealTied}(D)\Longrightarrow\texttt{VolTie}(D)
 \Longrightarrow
 D.\mathrm{vol}(W)=\operatorname{massI}(D,W.\mathrm{rank})
                                                                  \tag{RSL.5}
\]
for every supplied constant-signature $W$ satisfying
\texttt{SigOK}.  The first implication simply specializes the
\texttt{vol} field of \texttt{RealTied} to the constant signature
function.  Its measure and fundamental-domain quantifiers are
unchanged.  For the second, GV9--GV10 give the actual normalized
domain measure $M_n$, and RSL1 applied to the standard positive
space gives $M_n=\operatorname{massI}(D,n)$.  The actual
\texttt{VolTie} identifies that measure with
$\operatorname{ofReal}(D.\mathrm{vol}(W))$; positivity again removes
\texttt{ofReal}.  This proves the equality on the same $D$.

Apply this to the original $D.\mathrm{sp}.Y(0)$.
Its stored signature and rank identities give rank $D.R.r$, so
(RSL.5) is the literal \texttt{BaseVolumeEqualsMass D}.  Alternatively
one can use RSL1 in rank two, the actual finite class decomposition
of GV11, and the successive complements of RSL2, exactly as in
\texttt{volume\_eq\_mass\_rank\_two}.  Both routes retain the
original base space and its volume.

If a family has an eventual \texttt{RealTied} witness at threshold
$a$, this same $a$ supplies its eventual \texttt{VolTie} and
base-volume identities, with the already stipulated condition
$x>\mathcal F.x_1$ wherever the family predicate requests it.
The standard adjacent genus-mass ratio is separately supplied by
MF19.  Thus the fields \texttt{h2b} and \texttt{h2v} of
\texttt{FamGLmin} have no remaining same-object transport gap.
This does not produce \texttt{RealTied} for an arbitrary record:
DC20's normalization hypotheses, or an actual selected
$T.\mathrm{theta}.\mathrm{real}$, must still supply that record's
real-place witness.
% END RSL3

\subsection{Finite characters on the original coefficient field}

% BEGIN RSL4
\paragraph{RSL4: the required roots already belong to the chosen field.}
\label{claim:RSL4}
Fix the actual true-column data
\[
 I:\texttt{TC}(B,X),\quad I.\texttt{Ok},\quad
 l:I.\texttt{Laws},\quad Z:\texttt{Later}(B,I.\mathrm{cols}(l)),
 \qquad D=I.\mathrm{rec7}(l,Z).
                                                                  \tag{RSL.6}
\]
For $i=0,1$, put $p=g_i=D.\mathrm{host}.g(i)$.
The existing nontrivial additive character in $I.\texttt{Ok}$
supplies a primitive $p$th root of unity in the original $K$.
Here is the argument, which uses no coefficient-field enlargement.
The selected places over $g_i$ form a nonempty set: their cardinality
is $f=2\,\mathrm{ord}>0$.  At one of these places the original
residue map sends $g_i=\pi_i\overline\pi_i$ to zero.  As $g_i$ is
prime, the residue field has characteristic $p$.  The character
\[
                    \psi:(k_w,+)\longrightarrow K^\times
\]
stored in the character record is nontrivial, so choose $a$ with
$\zeta=\psi(a)\ne1$.  Additivity gives
$\zeta^p=\psi(pa)=\psi(0)=1$.  Its order is a nontrivial divisor
of the prime $p$, hence is exactly $p$.

Under the definition of \texttt{rec7}, this is literally the field
$D.\texttt{what}_2.\texttt{LL}$ and literally the same prime $g_i$.
The source also constructs roots of order $g_i-1$, but none is
needed in the zero-level argument below: there are no multiplicative
blocks at that level.  Keeping this distinction avoids imposing an
unnecessary cyclotomic extension or a stronger character condition.
% END RSL4

% BEGIN RSL5
\paragraph{RSL5: one explicit exponent works for any supplied symmetric forms.}
\label{claim:RSL5}
Suppose the original coefficient field contains primitive
$g_0$th and $g_1$th roots.  For every word $u$ and every actual
\texttt{FormsTiedSym} object $T$ on $D,u$, its level-zero character
satisfies
\[
 \chi_T(\kappa)^{g_0g_1}=1,
 \qquad \chi_T(\kappa)\in s_K(K),
 \qquad \kappa\in T.\mathrm{cls}(0).\mathrm{lev}.U.
                                                                  \tag{RSL.7}
\]
In particular this proves the exact predicate \texttt{ChiFin0 T},
with the explicit positive exponent $g_0g_1$.

To prove it, use the \texttt{glev0} field of this very $T$.
It expresses $\chi_T(\kappa)$ as the finite product, over the two
auxiliary primes and their selected places, of the inverses of the
local unipotent characters.  At level zero their multiplicative
block count is zero.  Each local term is therefore the value of
an additive character $\psi$ in residue characteristic $g_i$.
For every residue element $a$,
$\psi(a)^{g_i}=\psi(g_i a)=1$; the same holds for its inverse.
Each local term consequently has order dividing $g_0g_1$, and so
does their finite product.

For rationality, let $\zeta_i\in K$ be a primitive $g_i$th root.
The numbers $1,s_K(\zeta_i),\ldots,s_K(\zeta_i)^{g_i-1}$ are
distinct and are roots of $X^{g_i}-1$.  A polynomial of degree
$g_i$ has at most $g_i$ roots, so every complex $g_i$th root is
one of these powers.  Thus every local character value and its
inverse lie in $s_K(K)$; multiply their preimages in $K$ to obtain
the required preimage of $\chi_T(\kappa)$.  This proves both
clauses of (RSL.7), with no restriction on the particular supplied
$T$ beyond its original fields.

RSL4 gives the needed roots for (RSL.6).  TG12 supplies
$C:\texttt{ColsTiedSym}(I.\mathrm{rec7}(l,Z))$, with its original
standard-column and standard-level ties, from the supplied
$I.\texttt{Ok}$, pins, laws and later-stage record; DR3 supplies
the initial realizer used there.  Applying (RSL.7) to each
$C.\mathrm{forms}(u)$ proves the complete finite-character
endpoint in the source.
More strongly, (RSL.7) applies to any previously selected
$C$ on this same $D$.  Thus the older FC19 terminal-geometry
boundary is superseded in this scope by TG12, and we need not
replace a selector's columns by newly chosen ones.
% END RSL5

% BEGIN RSL6
\paragraph{RSL6: the same-selector and family quantifiers.}
\label{claim:RSL6}
Fix an original family $\mathcal F$ and a threshold selection
$s:\texttt{TSelF}(\mathcal F)$, with threshold $a_s$.
For $x>\mathcal F.x_1$, $a_s\le x$ and $i\in\mathcal F.I_x$,
write
\[
       T_{x,i}=s(x,x>\mathcal F.x_1,a_s\le x,i).T.
\]
Assume that at a threshold $a_\chi$ the roots of RSL4 have been
identified in the original field of $\mathcal F.D(x,i)$, for
every such $x,i$.  For example this follows when those same
records are the true-column \texttt{rec7} records of (RSL.6),
with an equality of records or an explicit transport of the
entire dependent tuple.  The conclusion is
\[
             \texttt{RealPlaceSecondLayerConditions}(\mathcal F).
                                                                  \tag{RSL.8}
\]
Take
\[
             a=\max\{a_\chi,a_s,\mathcal F.x_1+1\}.
                                                                  \tag{RSL.9}
\]
For every real $x\ge a$ and every $i\in\mathcal F.I_x$, retain
the original pair
\[
       T_{x,i}.\mathrm{theta}.\mathrm{real},\qquad
       T_{x,i}.\mathrm{cols}.
                                                                  \tag{RSL.10}
\]
RSL5 supplies \texttt{ChiFin0} for every word of these very
columns.  Equations (RSL.9)--(RSL.10) are exactly the witnesses
of (RSL.8), whose threshold precedes both $i$ and $u$.
The choice of the exponent $g_0g_1$ may depend on the family
member, as the predicate allows; no uniform field, prime or
character exponent is being assumed across members.

There is also a direct bundle formulation.  In
\texttt{WeilUnitaryEntries.ptE\_of\_wit}, a witness $q$ supplies
the real-place field $q.x.\mathrm{real}$ and the true-column
data $q.w.I,q.w.\mathrm{ok},q.w.\mathrm{pins},q.w.l,q.Z$.
RSL5 supplies its columns on precisely $q.w.I.\mathrm{rec7}$.
The linked-bundle equality transports the whole dependent tuple
$(D,O,\mathrm{Bundle})$, so it transports their conjunction to
the original $\mathcal F.D(x,i)$.  If these witnesses exist for
every $x>\mathcal F.x_1$ and every $i$, threshold
$\mathcal F.x_1+1$ suffices.  An equality of ranks, an equality
of abstract fields, or separately chosen real and column models
would not justify this step.

Finally, applying RSL3 member by member gives eventual
\texttt{VolTie}.  Applying RSL7 below to
$(T_{x,i}.\mathrm{cols}.\mathrm{forms}(u)).\mathrm{to0}$ gives
the actual trivial-stabilizer condition.  Together with RSL5
this supplies \texttt{tCVTies} and
\texttt{FamU ConcreteFormsLevelZero}; threshold $a$ already
forces $x>\mathcal F.x_1$.  Thus no additional eventual
stabilizer theorem is needed on this selected route.  For a
class of families, the outer quantifiers are those of
\texttt{Closed}: there is $\varepsilon_1>0$ such that for
$0<\varepsilon\le\varepsilon_1$ there is a transported bound
$\mathcal K$ valid for every $\mathcal F,C_H$ in the class
with $\mathcal F.\mathrm{Large}(\mathcal K,C_H)$.
The maximum in (RSL.9) is taken inside those family quantifiers,
and before every member index and word.  Thus the proof lifts
an already closed supplier without exchanging its quantifiers.
% END RSL6

\subsection{The actual integral and metrized induction model}

For this subsection fix $D,u$ and an actual
$T_0:\texttt{Forms0}(D,u)$.  In the preceding construction $T_0$
is the literal \texttt{to0} of the supplied symmetric forms.  Write
\[
 \begin{gathered}
 B=(D.\mathrm{fm}(u)).\mathrm{lat}(0),\qquad
 \mathscr C=(D.\mathrm{mem}(u)).C.\mathrm{cls}(0),\\
 \mathscr G=\prod_{v\in A.v_T} k_v^\times,
 \qquad A=(D.\mathrm{mem}(u)).A.
 \end{gathered}
                                                                  \tag{RSL.11}
\]
Thus $B$ is a Hermitian $\mathcal O_K$-lattice in its original
$K$-space $B.V$, with its original basis and metrics.  Each $k_v$
has cardinality $T=A.T$ and $|A.v_T|=f=D.R.f$.  The form-valued
realization $j:B.V\longrightarrow\{
\mathscr C\to\hbox{polynomial functions}\}$ is the original
semilinear embedding $(D.\mathrm{fm}(u)).\mathrm{emb}(0)$.
The fields \texttt{val0.rat}, \texttt{val0.lat} and
\texttt{norm0} in $T_0$ identify rational forms, the integral
lattice and the metrics, respectively.  None of these is inferred
merely from a dimension count.

% BEGIN RSL7
\paragraph{RSL7: two auxiliary primes kill every actual stabilizer.}
\label{claim:RSL7}
Every stabilizer of an original class $c\in\mathscr C$ is trivial.
This is the exact predicate \texttt{admS D u}; it follows already
from $T_0$, without a further period-admissibility or asymptotic
size hypothesis.

First consider a rational unitary matrix $\gamma$ such that
$y^{-1}\gamma y\in U$, where $y=T_0.\mathrm{cls0}.y(c)$ and
$U=T_0.\mathrm{cls0}.\mathrm{lev}.U$.  The level conditions imply
that this conjugate preserves the prescribed local lattices at
every finite place.  Their intersection is a full global
$\mathcal O_E$-lattice by LP1, and $\gamma$ preserves it.  The
level-zero Hermitian space in this construction is totally
positive definite at all archimedean embeddings.  Hence its
rational unitary lattice automorphism group is finite: after
restriction of scalars, the positive trace form bounds the images
of a fixed integral basis; a discrete lattice has only finitely
many points in each of these bounded balls.  In particular
$\gamma$ has finite order.

At either auxiliary prime $p=g_i$, the original
\texttt{glev0} conditions say that in the supplied split frame
the local matrix is integral and has upper-unipotent reduction.
Here the level-zero rank is $D.R.k$; the exact
\texttt{rank\_eq\_first\_level} identity gives this, so no
lower block with unrestricted eigenvalues remains.  The kernel
of reduction in $\mathrm{GL}_k$ over a complete discrete
valuation ring has no nonidentity torsion of order prime to
$p$.  To check this directly, suppose $M^n=1$, $p\nmid n$ and
$M=1+\varpi^a C$ with $a\ge1$ minimal and $C\not\equiv0$
modulo $\varpi$.  The binomial expansion modulo
$\varpi^{a+1}$ is
\[
                    M^n\equiv1+n\varpi^a C,
\]
which is not one, a contradiction.  If $M=1$ to every order,
separatedness of the complete ring gives $M=1$ directly.
An upper-unipotent matrix over a finite field of characteristic
$p$ has $p$-power order, since sufficiently large $p$-powers
kill its nilpotent part.  Combining these two facts shows that
every finite-order lift of such a matrix has $p$-power order.

Apply this at both $g_0$ and $g_1$.  The rational matrix maps
injectively into each split local factor, and conjugation does
not change its order.  That order is therefore both a
$g_0$-power and a $g_1$-power.  The distinctness field of the
original host gives $g_0\ne g_1$, so the order is one.
Finally \texttt{cls0.stabI} injects the abstract stabilizer
stored in $D$ into precisely this rational group, with exactly
this range.  Thus it too is the one-element group.  The class
weights in the original norm formula are consequently all one.
% END RSL7

% BEGIN RSL8
\paragraph{RSL8: actual diamonds and a free action on the actual classes.}
\label{claim:RSL8}
The group $\mathscr G$ in (RSL.11) acts freely on $\mathscr C$
by the original diamond dilations.  We give the local existence
and global freeness checks.

First the selected $T$-places are disjoint from the auxiliary
places.  The original \texttt{atT.piU\_cong} gives
$\pi\equiv1$ modulo
$N=Q_Sqd(g_0g_1)^2$, and \texttt{piU\_norm} gives
$\pi\overline\pi=T$.  At any prime above $g_i$, both
$\pi$ and $\overline\pi$ reduce to one, so $T$ reduces
to one.  It cannot equal $g_i$, whose residue is zero.
This derives $T\ne g_0,g_1$ from the supplied fields,
rather than treating it as an additional host condition.
The original tower-divisibility field also gives
$4\mid Q_S$, hence $4\mid N$.  The same norm congruence
therefore gives $T\equiv1\pmod4$.  Since $T$ is prime,
it follows in particular that $T\ge3$, pointwise and
without an asymptotic size argument.

At a selected place over $T$, the element $\pi$ in
$T_0.\mathrm{atT}$ has residue zero and $\overline\pi$ has
nonzero residue.  Put $\alpha=\pi-\overline\pi$; then
$\overline\alpha=-\alpha$ and $a=\alpha^2$ lies in $F$.
Choose an integral $b\in F$ lifting the nonzero residue of
$\alpha$ at this place.  The residue field is the prime field
of order $T$, and $T\ne2$, so
$b^2-a\in\mathfrak p$ and $2b$ is a unit.  Successive Newton
corrections for $X^2-a$ converge in the completed discrete
valuation ring to a unit $r$ with $r^2=a$.  The explicit
idempotents
\[
            e=\frac{1+\alpha/r}{2},\qquad
            1-e=\overline e
                                                                  \tag{RSL.12}
\]
are integral in the completed CM algebra and split its two
factors.

The level functional has an additive quotient of prime order
$T$.  An idempotent acts on this quotient by zero or identity:
its image is an additive subgroup, hence is zero or the whole
group, and a surjective idempotent is the identity.  Replacing
$e$ by $1-e$ if needed, take the one acting as identity.  For
$a_v\in k_v^\times$, choose an integer $k$ representing it,
with $T\nmid k$.  The local scalar
\[
                   \mu=k^{-1}e+k(1-e)
                                                                  \tag{RSL.13}
\]
is an integral unit of norm one.  It preserves the original
local lattices and scales the level functional by $k^{-1}$,
which is the prescribed diamond operation.  Put this scalar
at that place and the identity elsewhere.  This gives a
genuine finite adelic unitary element, not just a residue-class
operator.  It normalizes $U$, because scaling a level functional
by a unit preserves its kernel.  Elements at different places
commute, and its $(T-1)$st power belongs to $U$ by
$a_v^{T-1}=1$.  The quotient action is independent of the
integer lift.  The \texttt{fv0} and level fields of $T_0$
identify this action with the stored dilation of the original
class set and with the stored operators on forms.

To prove freeness, let a multidiamond $\delta$ representing
$g\in\mathscr G$ fix the class of $y$.  Thus
$y=\gamma y\delta\kappa$ for a rational unitary $\gamma$
and $\kappa\in U$, with the harmless inverse convention
absorbed into $g$.  Since $\delta$ normalizes $U$ and
$\delta^{T-1}\in U$, iteration gives
$y^{-1}\gamma^{T-1}y\in U$.  RSL7 implies
$\gamma^{T-1}=1$.  At the auxiliary primes $g_0,g_1$ the
diamond is the identity.  Therefore $y^{-1}\gamma y$ there
has the same upper-unipotent reduction as a level element.
The two-prime argument in RSL7 applies to this finite-order
$\gamma$ and gives $\gamma=1$.  It follows that
$\delta\in U$.  At each selected place the level functional
is nonzero on its quotient of order $T$; invariance under
its multiplier forces that multiplier to be one.  Thus every
coordinate of $g$ is one, which proves freeness.

All local fields, idempotents and functional quotients in this
argument belong to $T_0.\mathrm{atT}$ and
$T_0.\mathrm{cls0}.\mathrm{lev}$.  In particular the action is
the original one occurring in the definitions of
$\mathrm{ramSub}$ and \texttt{FreeT}.
% END RSL8

% BEGIN RSL9
\paragraph{RSL9: exact integral orbit slices and isometry at every embedding.}
\label{claim:RSL9}
Assume in addition \texttt{ChiFin0} for $T_0$.  There exist a
nonzero good Hermitian $\mathcal O_K$-lattice $W$ and an
injective $K$-linear map
\[
 \Psi:W.V^{\mathscr G}\longrightarrow B.V
                                                                  \tag{RSL.14}
\]
such that
\[
 \begin{split}
 B.L&=\Psi\bigl((W.L)^{\mathscr G}\bigr),\\
 \langle\Psi F,\Psi F'\rangle_\sigma
   &=\sum_{g\in\mathscr G}\langle F_g,F'_g\rangle_{W,\sigma}
       \quad\hbox{at every infinite place }\sigma,\\
 (D.\mathrm{fm}(u)).\mathrm{ramSub}
   &=\Psi\left(\bigcap_{v\in A.v_T}\ker\mathrm{Nm}_v\right),
 \end{split}                                                     \tag{RSL.15}
\]
Here $\mathrm{Nm}_v$ sums along the $k_v^\times$ coordinate.
The source calls these three properties
\texttt{InducedLattice}.

\emph{Rationality and the slice decomposition.}
Each original diamond operator preserves $B.L$, by the original
\texttt{rhoT\_lat} field, and therefore preserves $B.V$ because
$B.L$ spans it.  Injectivity of $j$ descends it uniquely to a
$K$-linear automorphism $\rho(g)$ of $B.V$; the group relations
descend by the same injectivity.  Thus the group action is on
the original rational space, not just its complexification.

Choose one representative in each of the finite free orbits of
RSL8, and let $W.V\subset B.V$ consist of forms supported on
these representatives.  Restricting a rational form to any set
of classes is still rational.  Indeed \texttt{val0.rat}
characterizes the image of $j$ by classwise rational polynomial
representatives; retain that polynomial on the chosen classes
and take zero elsewhere.  Trivial stabilizers, proved in RSL7,
make these classwise data legitimate automorphic forms, because
the only consistency condition inside a rational double coset
is its stabilizer condition.  Extending by the level character
and rational unitary action gives the required form, and its
classwise rationality puts it back in the image of $j$.

Define
\[
                       \Psi(F)=\sum_{g\in\mathscr G}\rho(g)F_g.
                                                                  \tag{RSL.16}
\]
At any class $c$, precisely one summand can be nonzero: freeness
gives exactly one $g$ carrying the appropriate class to the
chosen representative of its orbit.  Thus (RSL.16) is injective.
For a given form $b$, restrict it to each of the translated
representative sets and apply the inverse translate; these
pieces belong to $W.V$ and sum back to $b$.  This proves
surjectivity as well and describes the inverse explicitly.

\emph{Equality of the integral lattices.}
Put $W.L=B.L\cap W.V$.  The \texttt{val0.lat} field says that a
rational form lies in $B.L$ exactly when, at every class and
every coefficient place, its polynomial value lies in the
specified local integral polynomial module, in every allowed
local basis.  This is a classwise condition; zero satisfies it.
Consequently restriction to a set of classes preserves $B.L$.
Every $\rho(g)$ and $\rho(g^{-1})$ also preserves $B.L$.
The explicit inverse to (RSL.16) therefore preserves the
lattice coordinate by coordinate, and proves
\[
       \Psi(F)\in B.L\quad\Longleftrightarrow\quad
                    F_g\in W.L\ \hbox{for every }g.
                                                                  \tag{RSL.17}
\]
This is equality of integral modules, rather than commensurability.
The slice projector is $K$-linear and preserves $B.L$, so its
image $W.L$ is finitely generated and spans $W.V$; restricting
the positive definite metrics makes it a good Hermitian lattice.

It is nonzero.  The original class set contains the identity
double coset.  The level-zero coefficient representation has a
nonzero highest-weight function.  Explicitly, for a dominant
integer weight $\lambda_1\ge\cdots\ge\lambda_k$, take
\[
 (\det Z)^{\lambda_k}
       \prod_{j=1}^{k-1}\Delta_j(Z)^{\lambda_j-\lambda_{j+1}},
\]
where $\Delta_j$ is the leading $j\times j$ minor, and take
the product over the stored embeddings.  The minor exponents
are nonnegative, the determinant factor is defined on the
invertible frames even when $\lambda_k<0$, and its value on
the coordinate identity frame is one.  This is the extreme
weight element of the original coefficient module, expressed
in the designated frame coordinates.  It therefore gives a
nonzero value function of exactly the required weight.
Placing it on one class and zero elsewhere gives a nonzero
complex form, since the stabilizer is trivial.  The original
\texttt{emb\_basis} spans this complex form space; hence
$B.V\ne0$.  Since $B.L$ spans $B.V$ and (RSL.16) is a
bijection, $W.V$ and $W.L$ are nonzero as claimed.

\emph{The original metrics.}
The field \texttt{norm0} expresses each squared norm as the
sum over classes of the specified Fock norms of its polynomial
values, weighted by reciprocal stabilizer orders.  RSL7 makes
these weights one.  Forms supported on disjoint sets of classes
are orthogonal at every embedding.  One can check the mixed
term from the norm identity: apply it to $x+y$ and $x+a y$,
where $a\in K$ has nonreal image at the embedding in question.
Such an $a$ exists because $K$ is totally complex.  The two
equalities force both real and imaginary parts of the Hermitian
cross term to be zero.

It remains to prove that each $\rho(g)$ is an isometry; a
permutation of supports alone would not suffice.  For a single
diamond and each class, the original double-coset equality and
\texttt{val0} express its new polynomial value as the old value
after substitution by a rational unitary matrix $\gamma^{-1}$,
multiplied by a level-character value $\chi(\kappa)$, and
evaluated at the dilated class.  The character hypothesis gives
$r\in K$ and $N>0$ with
$s_K(r)=\chi(\kappa)$ and $s_K(r)^N=1$.
Injectivity gives $r^N=1$ in $K$, so
\[
                         |\sigma(r)|=1
                   \quad\hbox{for every embedding }\sigma.
                                                                  \tag{RSL.18}
\]
The rational unitary matrix has its entries in the original
$E$.  Mapping its entries to $K$ through the stored embeddings
commutes coefficient by coefficient with polynomial
substitution and with every subsequent embedding $\sigma$.
The source value function also contains its determinant shift:
the corresponding polynomial transport has the factor
$\prod_\tau\det(\gamma_\tau)^{\mathrm{shift}_\tau}$.
Each $\gamma_\tau$ is unitary for its Hermitian form, so
$|\det(\gamma_\tau)|=1$ at every embedding.  This factor
therefore changes no squared norm, including for negative
integer shifts.

The same unitary identity carries an orthonormal frame at
$\sigma$ to another orthonormal frame.  Apply \texttt{norm0}
to the rational restriction supported on just one class,
whose existence was proved above.  Its sum has only that
class's contribution, so this contribution is independent of
the chosen orthonormal frame.  This justifies changing frames
separately for each class even though $\gamma$ can depend on
that class.  Changing frame thus preserves the evaluated
polynomial norm, while (RSL.18) preserves the character
factor.  Summing over the permuted classes
shows that the original squared norm is invariant under the
diamond at every $\sigma$.  Polarization as above gives
invariance of the full Hermitian pairing.  Products of diamonds
are therefore isometries too.

Distinct translates of the slice have disjoint supports.  The
preceding two paragraphs now prove the pairing identity in
(RSL.15), with the restricted metric on $W$.

\emph{The actual ramified subspace.}
For a single coordinate $v$ and $a\in k_v^\times$, (RSL.16)
intertwines $\rho(a)$ with translation of that coordinate of
the function $F$.  Summing over $a$ intertwines the original
norm operator with coordinate summation.  The definition of
the original $\mathrm{ramSub}$ is the joint kernel of exactly
these norm operators; the original rational embedding is
injective.  This proves the last equality of (RSL.15).
Taking its intersection with (RSL.17) gives the exact original
ramified lattice, including its induced metric.
% END RSL9

% BEGIN RSL10
\paragraph{RSL10: the actual rank, mean and maximal-slope formulas.}
\label{claim:RSL10}
The source-facing conclusion retains
\texttt{HidaSetting D u}.  The argument below needs only
$T\ge3$, which is already pointwise true by RSL8.
Put $q=T-1$, and set
\[
 U=\mathcal O_K^{q},\qquad
 A_0=\ker\bigl(U\xrightarrow{\sum}\mathcal O_K\bigr).
\]
Give $U$ its standard metric.  RSL9 identifies the original
full lattice with $W\otimes U^{\otimes f}$ and the original
ramified lattice with $W\otimes A_0^{\otimes f}$,
integrally and isometrically.  The second identification uses
the integral splitting
$U=A_0\oplus\mathcal O_K e_q$: its tensor product shows that
the joint coordinate-sum kernel is precisely the indicated
tensor product, without a saturation or index correction.

If $h=\operatorname{rk}W$, their ranks are $h q^f$ and
$h(q-1)^f$.  Define the actual source integer
\[
 d_u=(D.\mathrm{mem}(u)).A.du
      =\dim_{\mathbb C}(D.\mathrm{mem}(u)).A.M_{\mathrm{ram}}.
\]
The original $\texttt{emb\_basis}$ identifies
$B.V\otimes_{K,s_K}\mathbb C$ with $A.M_{u,0}$,
so $D.duS(u)=h q^f$.  Each diamond and norm operator
descends to $K$, as proved in RSL9.  Their joint kernel
commutes with scalar extension: stack their finitely many
matrices into one matrix over $K$; its rank is unchanged
over $\mathbb C$, because rank is characterized by its
nonzero minors.  Under the original embedding, their
complex joint kernel is exactly $A.M_{\mathrm{ram}}$.
Therefore the original rational \texttt{ramSub} has
$K$-dimension $d_u=h(q-1)^f$.  The injectivity and range
of the original $D.\mathrm{nuIn}$ identify the division
vector space with this very rational kernel.  Thus
\[
 \frac{d_u}{D.duS(u)}
       =\left(1-\frac1{T-1}\right)^f.                 \tag{RSL.19}
\]
The lattice $W$ is nonzero and $q-1>0$, so both ranks are
positive.  The basis $e_j-e_q$ of $A_0$ has Gram matrix
$I+J$ and determinant $q$.  The Kronecker determinant
calculation in LC7, now applied through the exact isometry of
RSL9, gives
\[
 \boxed{\quad
 (D.\mathrm{fm}(u)).\mathrm{ramLat}.\mu
       =B.\mu-\frac{f\log(T-1)}{T-2}.
 \quad}                                                          \tag{RSL.20}
\]
For clarity, the tensor Gram determinant is
$q^{f(q-1)^{f-1}}$ when $f>0$; dividing its logarithm by
$(q-1)^f$ gives the loss in (RSL.20).  Finite indices in a
free sublattice of the projective $W.L$ occur with the same
normalized multiplicity on both sides, so no freeness
assumption on $W.L$ is required.  The squared-metric slope
convention gives exactly (RSL.20), with no extra factor two.
When $f=0$ the tensor product is $W$, the rank factor is one
and the loss is zero, so this boundary case is consistent too.

The original ramified lattice is the intersection of $B.L$
with a rational subspace, with the restricted original
metrics.  Every lattice in this subspace used to compute its
maximal slope is also an admissible lattice in $B$.  Taking
the supremum gives
\[
        \mu_{\max}\bigl((D.\mathrm{fm}(u)).\mathrm{ramLat}\bigr)
                           \le\mu_{\max}(B).                     \tag{RSL.21}
\]
This is also valid if the convention restricts to saturated
submodules: the ramified intersection is saturated in $B.L$,
and saturation inside it remains saturated in $B.L$.

Finally choose a complex basis of $W.V\otimes_{K,s_K}\mathbb C$.
Through (RSL.16), the full form space becomes functions on the
finite free $\mathscr G$-set consisting of one copy of
$\mathscr G$ for each basis vector.  The single-coordinate
actions intertwine the stored $\rho_T$ operators.  This is
the literal \texttt{FreeT D u} record.  Equation (RSL.20) is
the literal \texttt{SlopeLine D u}; (RSL.19)--(RSL.21) are
the three fields of \texttt{family\_property\_and} on the
original data.  RSL8 alone already gives the corresponding
complex free model, without the character hypothesis; the
finite-character condition is used to obtain the isometry at
every coefficient embedding.

Applied to RSL6 at one common threshold, these results supply
the fields \texttt{hW} and $h_\mu$ of
\texttt{FamGLmin}.  This closes LC7's actual integral and
isometric transport boundary.  It supplies no assertion
about the Hecke maps between the two different local
integral lattices appearing in RSL13.
% END RSL10

\subsection{A smaller fullness cut and the residual global ledger}

% BEGIN RSL11
\paragraph{RSL11: actual fullness needs only the ramified local inclusion.}
\label{claim:RSL11}
For an original family $\mathcal F$, suppose there is one
threshold $a_N$ such that, for every $x\ge a_N$, every
$i\in\mathcal F.I_x$, every word $u$ and every coefficient
place $w$ of residue characteristic $T_u$,
\[
                N_\lambda(D,u,w)\subset S_{\mathrm{ram}}(D,u,w),
       \qquad D=\mathcal F.D(x,i).                    \tag{RSL.22}
\]
Then \texttt{NuFullF Fm} follows from the original fields of
\texttt{FamP}.  In particular it does not require the entire
congruence length bound, a mean-slope theorem, or a new
arithmetic existence assertion.

Fix a member beyond $\max\{a_N,\mathcal F.x_1+1\}$, and write
$M=L.\mathrm{NuLat}(u)$ and
$R=\mathrm{MramL}(D,u)$, both in the original vector space
$(D.\mathrm{div}.\mathrm{Nu}(u)).V$.  The field
$(\mathcal F.\mathrm{ex}(x,i)).\mathrm{lat}.\mathrm{atT}$
identifies the localization of $M$ at $w\mid T_u$, through
the original $D.\mathrm{nuIn}(u)$, with $N_\lambda(D,u,w)$.
The first component of $S_{\mathrm{ram}}$ is the original
local full lattice $S(0,w)$.  The field $D.\mathrm{locSpan}$
clears its local denominators.  Thus (RSL.22) implies
$M_w\subset R_w$ at every $w\mid T_u$.
At the other places, the original
$\mathcal F.\mathrm{nuDef}(x,i)$ supplies this inclusion.
Consequently $M\subset R$ globally.  For completeness, if
$z\in M$ did not belong to $R$, its annihilator ideal
$\{a\in\mathcal O_K:az\in R\}$ would be proper and
contained in a maximal ideal $\mathfrak p$.  Local inclusion
at $\mathfrak p$ would supply $s\notin\mathfrak p$ with
$sz\in R$, a contradiction.

The original $R$ is finitely generated: its injective map
$D.\mathrm{nuIn}$ identifies it with a submodule of the
finitely generated lattice $B.L$, and $\mathcal O_K$ is
Noetherian.  Since
$\mathcal O_K$ is Noetherian, its submodule $M$ is finitely
generated too.  The original \texttt{NuDef} also gives
$M_w=R_w$ at a suitable place away from $T_u$.  There is
always an eligible rational prime: use $2$ when $T_u\ne2$,
because $2$ belongs to the host's designated small-prime set;
when $T_u=2$, use the host prime $q$, which is distinct from
$2$ and has the designated non-inert class.  Choose a
coefficient place above this prime.  For finitely many
generators $r_1,\ldots,r_b$ of $R$, local equality gives
$s_j\notin\mathfrak p$ with $s_jr_j\in M$.
Their product $c=\prod_j s_j$ is nonzero and satisfies
\[
                              cR\subset M\subset R.             \tag{RSL.23}
\]
This is exactly \texttt{NuCommens D L}.

The original $R$ spans the original division vector space:
apply $D.\mathrm{nuIn}$ to a vector, express it in the
$K$-span of $B.L$, clear the finitely many coefficients,
and pull back.  Thus (RSL.23) proves that $M$ spans it too.
Finally the actual \texttt{nprimeDef} field of the same
$\mathcal F$ gives
\[
                 N=(D.\mathrm{div}.\mathrm{Nu}(u)).L=I_u M,
                                                                  \tag{RSL.24}
\]
with the same nonzero fractional ideal $I_u$ as in MT1.
Choose $0\ne d\in\mathcal O_K$ with $dI_u\subset\mathcal O_K$.
Then $dN\subset M$, so Noetherianity and injectivity of
multiplication by $d$ imply finite generation of $N$.
Choose a nonzero integral element $a\in I_u$; (RSL.23)
gives $acR\subset N$, so $N$ spans the division vector
space.  These are both fields of $D.\texttt{NuFull}$.
The threshold was chosen before $i,u,w$, hence proves
\texttt{NuFullF} with its original uniformity.  Neither
\texttt{nuDef} nor \texttt{nprimeDef} has been added as a
new datum; both are fields of the original family.
% END RSL11

% BEGIN RSL12
\paragraph{RSL12: actual positive rank and the sufficient operator cut.}
\label{claim:RSL12}
Suppose an original family has the actual second layer of
RSL6, \texttt{TuSize}, and the eventual inclusion (RSL.22).
Then
\[
                   \texttt{NuFullF}\ \wedge\ \texttt{J.MuRank}
                                                                  \tag{RSL.25}
\]
holds for that same family.  The second conjunct says that
each original division lattice has strictly positive rank.

Indeed RSL8 gives $T_u\equiv1\pmod4$ and hence $T_u>2$
on each of the original records.  Take one maximum of
the second-layer threshold, $a_N$, and
$\mathcal F.x_1+1$.  RSL11 gives fullness of the original
division lattice.  Its full $K$-rank equals the dimension
$d_u=(D.\mathrm{mem}(u)).A.du$ of the original ramified form space, since
$D.\mathrm{nuIn}$ is injective with precisely that range.
RSL9--RSL10 give $D.duS(u)>0$, $T_u>2$ and (RSL.19),
so $d_u>0$.  This proves (RSL.25) at one threshold.
The nonzero slice constructed in RSL9 makes the positivity
argument slightly stronger than the generic source lemma
\texttt{muRank\_min}, which instead obtains $D.duS(u)>0$
from \texttt{ChosenProviso}.  On the selected final route
\texttt{ChosenE Fm 3 CH} already contains that proviso;
either argument applies, and neither changes the family.
The retained \texttt{TuSize} hypothesis is consequently
unused in this stronger actual-forms proof.  In particular
we have not appealed here to the longer numerical supplier
\texttt{TuSize}$\Rightarrow$\texttt{SettingEventually}.

We also isolate exactly what is needed to supply (RSL.22).
At a seen coefficient place above $T_u$, suppose the actual
level-zero operators satisfy
\[
 U_v(B_w(v))\subset S_w,\qquad
 U_v(M_{\mathrm{ram}})\subset M_{\mathrm{ram}},\qquad
 T_u^{b_v}S_w\subset B_w(v),
                                                                  \tag{RSL.26}
\]
where $b_v=\mathrm{bTop}(\tau_w(v))\ge0$ for the selected
places, as supplied by \texttt{HidaSetting}.  The normalized
operator $G_v=T_u^{b_v}U_v$ then preserves
$S_{\mathrm{ram},w}=S_w\cap M_{\mathrm{ram},w}$.
Rationality of $U_v$ follows from its integral map between
two full local lattices: choose a basis, clear denominators
in the full source lattice, and descend by the injective
original rational embedding.  The exact source definition
of $N_\lambda$ is the image of $S_{\mathrm{ram},w}$ under
the $m$th power of the product of these $G_v$, in the
source's ordered list of selected places.  Distributing the
central scalar factors gives the exponent
$m\sum_v b_v=\texttt{expT}$ in that original definition;
no interchange of two Hecke operators is used here.
Successive preservation proves (RSL.22).  At an unseen
place the selected set is empty, its product is the
identity, and $N_\lambda=S_{\mathrm{ram},w}$ directly.

Thus (RSL.26) is sufficient for fullness and positive rank.
The full three contracts are
\[
 \begin{gathered}
 \texttt{HeckeIntegralLevelZero},\qquad
 \texttt{HeckeRamifiedLevelZero},\\
 \texttt{SecondBranchInclusionLength}.
 \end{gathered}
\]
They imply (RSL.26), but
their reverse-operator identity and exact length formula
are stronger than this use requires.  Those stronger
parts are still needed in the selected global error
ledger.  We have proved the implication from (RSL.26),
not silently asserted those operator and length contracts.
% END RSL12

% BEGIN RSL13
\paragraph{RSL13: the exact global cut and the remaining proof tasks.}
\label{claim:RSL13}
For one fixed class $\mathcal C$ the selected source interface
is the eleven-field conjunction \texttt{FamGLmin C 4}.
The following ledger is literal; every field concerns the
same original family, member, coefficient field and word.
\begin{center}
\begin{tabular}{p{0.11\textwidth}p{0.47\textwidth}p{0.30\textwidth}}
\hline
Field & Contract & Human status in this chapter \\
\hline
$h_{2b}$ & Eventual standard adjacent genus-mass ratio & MF19, closed. \\
$h_{2v}$ & Eventual original base-volume equals mass & RSL1--RSL3, closed from the supplied real tie. \\
$h_3$ & \texttt{SigmaDegreeTwoSidedBoundOfGap} & B: explicit finite-product bounds and original-row installation; deep discrete-series producer covered by DC1, DC10, DC19. \\
$h_{Za}$ & \texttt{FamilySlopeLowerOpenPart} & B: exact column-lattice slope and standard-weight estimate. \\
$h_{Zb}$ & \texttt{FamilySlopeMaxOpenPart Fm 4} & B: original class-map degree and hook bound. \\
$tCV$ & Eventual \texttt{Nonempty (ColVal D u)} & B: actual polynomial/local-basis/archimedean dictionary. \\
$h_W$ & \texttt{FamU (Nonempty FreeT)} & RSL7--RSL10, closed from the second layer. \\
$h_\mu$ & \texttt{FamU SlopeLine} & RSL9--RSL10, closed from the second layer. \\
$h_I$ & \texttt{FamU HeckeIntegralLevelZero} & B: actual two integral Hecke maps. \\
$h_R$ & \texttt{FamU HeckeRamifiedLevelZero} & B: actual ramified stability and inverse composite. \\
$h_B$ & \texttt{FamU SecondBranchInclusionLength} & B: actual local inclusion and exact length. \\
\hline
\end{tabular}
\end{center}

There is a separate B installation gate when these conditional
\texttt{FamU} fields are used on actual late members:
\[
       \texttt{TuSize}\ \Longrightarrow\
       \texttt{SettingEventually},
                                                                  \tag{RSL.28}
\]
where the latter asserts the full \texttt{HidaSetting} at
one threshold before $i,u$.  The source supplier is
\texttt{Places.settingF\_of}, through
\texttt{numsL\_wtsL\_fam} and
\texttt{levelPrime\_setting\_of\_sizes}.
Its numerical size and weight comparisons still need a human
reconstruction.  In particular, the class projection supplying
\texttt{TuSize} does not itself prove all the alcove and
$b_{\rm top}\ge0$ inequalities.  The strong Hecke/length route
to (RSL.26), and its use in the final global ledger, must retain
this gate.  The weaker direct fullness input (RSL.22) does not
need it, and RSL12 has avoided it for positivity by the original
congruence $T\equiv1\pmod4$.

The smallest \emph{literal GL4 cut at this interface} consists
of the three fields $h_{Za},h_{Zb},tCV$.  Its sufficiency for
these three slots is by their definitions; no assertion of
absolute mathematical minimality is intended.  More concrete
producer targets are:
\begingroup\raggedright
\begin{itemize}
\item \texttt{column\_lattice\_exact\_slope}
in \path{GL4/ColumnLatticeExactSlope.lean}, followed by
the standard Schur-lattice slope and hook-length estimates
used in \path{GL4/LowerSlopeInequalityFamily.lean}.
The equality must hold for the actual \texttt{ColVal.Q}
with its exact base-change lattice and metric.
\item \texttt{class\_map\_archimedean}
in \path{GL4/ClassMapArchimedean.lean}, followed by the
class-map slope inequality and the constant-four hook bound
in \path{GL4/PStmtFree.lean}.  The local class bases and
the original archimedean matrices must be used together.
\item \texttt{ColValS.colValOfTiedS\_holds} in
\path{GlobalTheory/TransposeBaseMatrixInfinity.lean}.
Its useful smaller dictionary target is
\texttt{ColTieW}: an injective polynomial realization of
the original rational form space, the exact local integral
membership in bases of the actual $Q_{c,\tau}$, and the
exact sum-of-polynomial-norms formula for the original
metric.  RSL9 establishes the rational restriction and
metric transport mechanisms, but does not by itself
identify these $Q_{c,\tau}$ or write the complete
\texttt{ColTieW} constructor.
\end{itemize}
\endgroup

For the three GL6 operator slots, the exact source producers are:
\begingroup\raggedright
\begin{itemize}
\item \texttt{integral\_at\_t0\_lattice} in
\path{OrbitalBoundsConclusion/IntegralAtT0Lattice.lean};
\item \texttt{hecke\_ramified\_zero\_all} in
\path{GlobalTheory/UnramifiedTrace/HeckeRamifiedZeroAll.lean},
which installs the same-place matching and unramified
composite calculations;
\item \texttt{second\_branch\_all\_members} in
\path{GlobalTheory/CoordinateLengths/SecondBranchAllMembers.lean}.
\end{itemize}
\endgroup
Their mathematical content is a finite local coset/integrality
calculation and a discrete-valuation length calculation;
their exact installations remain B reconstruction tasks.
The finite product of operators in RSL12 proves their
weaker fullness consequence once (RSL.26) is supplied.
For the full ledger one must also write the finite-index
length estimates for the composite and its $m$th power;
the source declaration is
\texttt{congruence\_no\_kernel\_minimal}.

The closed transport claims do not reinstate older
boundaries.  TG12 supplies the full columns on true-column
data; DC15 and DC19 supply the large-gap discrete-series
existence, uniqueness and degree under their named
normalizations; DC20 matches all four real-place fields
when those normalizations are supplied.  FR1 and GF8
prove the finite-dimensional frame and normalized
Harder--Narasimhan consequences.  They do not identify
the actual GL4 lattices or the two Hecke lattices, and
they do not check by themselves every input of the final
\texttt{linkGLOBAL7\_items} assembly.  That last numerical
and lattice assembly remains an explicit B proof-of-implication
task, rather than being counted as closed because the
formal source calls it.

For precision, finite conjunction of already proved
\texttt{Closed} contracts is elementary.  Choose the
minimum of their finitely many positive $\varepsilon$
bounds.  For each allowed $\varepsilon$, take the supremum
of their finitely many transported bounds.  The definition
of \texttt{Large} at that supremum implies each original
\texttt{Large} condition.  Thus every contract holds for
one and the same $\mathcal F,C_H$; within that family,
take the maximum of the finitely many eventual thresholds.
This proves the finite joining step without changing
any outer quantifier.

On the selected final route put $\mathcal C=D3(3,\mathcal A)$.
The precise terminal source contract is
\[
 \begin{gathered}
 \texttt{Duty}\hbox{ on }\mathcal A,\quad
 \texttt{Closed}(\mathcal C,\texttt{RealPlaceSecondLayerConditions}),\\
 \forall D\,\texttt{MassRatioUnderVolume}(D),\quad
 \forall D\,\texttt{GenusMassAt}(D,2)\\
 \Longrightarrow\quad
 \texttt{Closed}\bigl(\mathcal C,
  \texttt{ChosenE}(3,C_H)\Rightarrow
                  \texttt{NuFullF}\wedge\texttt{MuRank}\bigr),\\
 \texttt{Closed}\bigl(\mathcal C,
   \texttt{TuSize}\Rightarrow\texttt{DGamSize}\Rightarrow
                  \texttt{LinkGLOBAL}(2,C_H,3,2,0)\bigr).
 \end{gathered}                                                   \tag{RSL.27}
\]
RSL1--RSL3 close both universal mass/volume inputs.
RSL4--RSL6 close the finite-character and real/column
join once the same-record selection has been installed.
The class projection $D3(3,\mathcal A)\to\mathcal A$
then supplies the original \texttt{ChosenWall} from
\texttt{Duty}; it does not choose a new family.
Membership in this same $D3$ class already includes
\texttt{TuSize}.  Consequently RSL12's \texttt{TuSize}
hypothesis is a class projection on the selected
fullness/rank route, not a new producer to construct.
The numerical parameters are exactly $C_{\rm abs}=2$,
$C_0=3$, $e_D=2$, $e_L=0$; $C_5=4$ belongs only to
the intermediate slope estimate.  The external
\texttt{DGamSize} argument is retained in (RSL.27)
although the \texttt{link8} wrapper does not use it.
To count the full human implication (RSL.27) as finished,
the B entries, the size-installation gate (RSL.28), and the
final assembly just identified
must still be written.  This chapter does not replace
those proofs by an invocation of \texttt{link\_gl}.
% END RSL13
